\documentclass[11pt]{article}

\usepackage[margin=1in]{geometry}
\usepackage{amsmath, amssymb, amsthm, mathtools}
\usepackage{enumitem}
\usepackage[colorlinks=true, linkcolor=blue, citecolor=blue,
            urlcolor=blue]{hyperref}
\usepackage{microtype}

% =====================================================================
% Notation
% =====================================================================
\newcommand{\HH}{\mathbb{H}}
\newcommand{\OO}{\mathbb{O}}
\newcommand{\Cl}{\mathrm{Cl}}
\newcommand{\Q}{\mathbb{Q}}
\newcommand{\R}{\mathbb{R}}
\newcommand{\Z}{\mathbb{Z}}
\newcommand{\spec}{\operatorname{spec}}
\newcommand{\op}{\mathrm{op}}
\newcommand{\Schur}{\operatorname{Schur}}
\newcommand{\id}{\operatorname{id}}
\DeclareMathOperator*{\Span}{span}

\newtheorem{theorem}{Theorem}[section]
\newtheorem{proposition}[theorem]{Proposition}
\newtheorem{lemma}[theorem]{Lemma}
\newtheorem{corollary}[theorem]{Corollary}
\theoremstyle{definition}
\newtheorem{definition}[theorem]{Definition}
\newtheorem{remark}[theorem]{Remark}
\newtheorem{notation}[theorem]{Notation}

% =====================================================================

\title{Cascade Selection on the $H_4$ Branch:\\
An Auxiliary Block Graph Laplacian and a Witnessed $(O1)$ Discharge
on $V_{600}$}

\author{Cascade--Classical Correspondence Programme}
\date{\today}

\begin{document}

\maketitle

\begin{abstract}
We give a finite-level witnessed $(O1)$ tuple for the
cascade-mechanism consumer-API obligation
\cite[\S 3.4]{CascadeMechanism} on the $H_4$ branch at every
finite top-rung level $N \geq 0$, parametrised by a free choice of
$\xi^{*} \in (F^0)^{W(H_4)} \setminus \{0\}$. The witness operator
is a new auxiliary block graph Laplacian
$L_N^{\mathrm{block}}$ on the finite-level vertex-sector Hilbert
space $X_N^{0, H_4}$ of \cite[Definition~\texttt{def:Xn-0}]{RefinementSpaces},
introduced in this paper and explicitly distinct from
\cite{ClosureDynamics}'s vertex / top-left block
$A_n^{\mathrm{vertex}} = d_n^{*} d_n$ (which acts trivially off
the $\mathrm{Im}(\OO)$ direction in $F^{0}$) and from
\cite{RefinementCompat}'s $\widetilde A_n$ (which lives in an
abstract scalar pure-midpoint reduction).

The paper proves five ingredients:
$L_N^{\mathrm{block}}$ is bounded self-adjoint and positive
semi-definite on $X_N^{0, H_4}$, with an explicit Chung-style
operator-norm bound (Proposition~\ref{prop:psd});
$G_N^{H_4}$ is connected at every level and
$\ker L_N^{\mathrm{block}} \cong F^0$ via constant-function
embedding ($32$-dimensional, Theorem~\ref{thm:kernel});
the gradient flow $e^{-t L_N^{\mathrm{block}}}$ converges to the
kernel projection with exponential rate governed by the
spectral gap $\lambda_{+}(L_N^{\mathrm{block}}) > 0$
(uniform over initial data at fixed $N$, Theorem~\ref{thm:convergence});
$\dim_{\R} (F^0)^{W(H_4)} = 28$ by irreducibility of the standard
$H_4$-reflection representation (Theorem~\ref{thm:fixed-dim});
and the resulting tuple satisfies $(O1)$ for the chosen $\xi^{*}$
and chosen top-rung initial datum (Theorem~\ref{thm:main}), with
the standing analytic preconditions $(O0)$ verified at the
top-rung level (Proposition~\ref{prop:o0}).

We do \emph{not} discharge $(O2)$, $(O3)$, the inverse limit, the
$D_4$ branch, or the strong sense of selection that would pick
a single ray in $(F^0)^{W(H_4)}$. The latter is named
\textbf{(H-uniqueness)} and lies outside the cascade-closure
consumer API as stated in
\cite[\S 3.4 lines 637--639]{CascadeMechanism}.
\end{abstract}

\tableofcontents

% =====================================================================
\section{Recap and scope}
\label{sec:intro}
% =====================================================================

\subsection{The narrow goal}

This paper is a single-target follow-up to
\cite{CascadeMechanism}: we supply a witnessed $(O1)$ tuple for
the cascade-closure consumer API on the $H_4$ branch at every
finite top-rung level $N \geq 0$. The
\cite[\S 3.4]{CascadeMechanism} consumer API consists of four
\emph{separate} obligations $(O0)$, $(O1)$, $(O2)$, $(O3)$
(\cite[lines 594--627]{CascadeMechanism}); we discharge $(O1)$
and verify the top-rung analytic pieces of $(O0)$ as a standing
precondition (the inter-rung pushforward bonding map needed by
the rung-sequence formulation of $(O0)$ is available from
\cite[Definition~\texttt{def:p0}, Theorem~\texttt{thm:bonding}]
{RefinementSpaces} but is not used in the present
$(O1)$-witness construction). $(O2)$ and $(O3)$
are not addressed and remain in their current
\cite{RefinementCompat} status (i.e., $(O2)$ \emph{as stated}
remains false; $(O3)$ holds in the abstract scalar pure-midpoint
model only).

The $(O1)$ discharge is \emph{parametrised} by a free choice of
$\xi^{*} \in (F^0)^{W(H_4)} \setminus \{0\}$ (the
$W(H_4)$-fixed subspace, of explicit real dimension $28$,
Theorem~\ref{thm:fixed-dim}). For each fixed $\xi^{*}$ we
exhibit an explicit five-tuple
$(\Phi_N, R_N, \Psi_N^{t}, \phi^{\star}_{N}, \phi_N^{(0)})$
satisfying every part of $(O1)$
(\cite[lines 609--617]{CascadeMechanism}). The freedom in
$\xi^{*}$ is the $(O1)$-witness parameter; the
\cite[\S 3.4 lines 637--639]{CascadeMechanism} statement
``selection of $\phi^{\star}_{N}$ itself is not provided by
cascade closure and is the successor paper's own task'' is
satisfied by our construction supplying \emph{a} $\phi^{\star}_{N}$.
Reducing to a single ray in $(F^0)^{W(H_4)} \setminus \{0\}$
requires the named open hypothesis \textbf{(H-uniqueness)}
(\S\ref{sec:open-hyp}), which lies \emph{outside} the
cascade-closure API.

\subsection{Explicit non-claims}

\begin{remark}[Non-claims box]
\label{rmk:non-claims}
This paper does \emph{not} establish any of the following:
\begin{enumerate}[label=(\roman*), leftmargin=2em]
\item \textbf{Single-state selection.} $\xi^{*}$ is freely chosen
  from $(F^0)^{W(H_4)} \setminus \{0\}$ (a $28$-dimensional cone);
  reducing to a single ray requires \textbf{(H-uniqueness)}
  (\S\ref{sec:open-hyp}). The cascade-closure consumer API does
  not provide single-ray selection.
\item \textbf{$(O2)$, $(O3)$, or the full $(O0)$--$(O3)$ API.}
  $(O0)$ is verified only at the top-rung level as a standing
  precondition for the witnessed $(O1)$ tuple (the full
  rung-sequence formulation of $(O0)$ is not built here; the
  required inter-rung pushforward bonding map is recorded as
  available from
  \cite[Definition~\texttt{def:p0}, Theorem~\texttt{thm:bonding}]
  {RefinementSpaces} but not used). $(O2)$ as stated remains false
  (\cite[Proposition~\texttt{prop:strict-fails}]{RefinementCompat}),
  and $(O3)$ holds only in the abstract scalar pure-midpoint
  model
  (\cite[Theorem~\texttt{thm:O3-discharge}]{RefinementCompat}).
  Neither is touched here.
\item \textbf{Inverse-limit $(O1)$.} We work strictly at finite
  top-rung level $N$; lifting to $N \to \infty$ would require
  the named open hypotheses of
  \cite[\S\texttt{sec:scope-and-gap}]{RefinementCompat}, namely
  $(L1)$ (boundary-vertex-incidence agreement of the level-$(n+1)$
  vertex Laplacian with the abstract scalar model up to a
  midpoint-supported defect) and $(L3)$ (energy-saturation of the
  $F^{0}$-valued harmonic-extension operator $H_{n+1}$). Neither is
  invoked here.
\item \textbf{$D_4$ branch.} The construction uses the
  $H_4$-branch action of \cite[Definition~\texttt{def:coxeter}]{RefinementSpaces}
  in a load-bearing way (clause 4); a $D_4$-branch analogue is
  a separate paper.
\item \textbf{Identification with upstream operators.}
  $L_N^{\mathrm{block}}$ is a new auxiliary operator introduced
  here. Theorems about it do not transfer to
  \cite{ClosureDynamics}'s vertex / top-left block
  $A_n^{\mathrm{vertex}} = d_n^{*} d_n$ or
  \cite{RefinementCompat}'s $\widetilde A_n$.
\item \textbf{$W(H_4)$-non-equivariant initial data.}
  Convergence still holds (Theorem~\ref{thm:convergence}), but
  the limit lies in the full $32$-dimensional
  $\ker L_N^{\mathrm{block}}$, not in the $28$-dimensional
  $\iota_N\bigl((F^0)^{W(H_4)}\bigr)$.
\item \textbf{Continuum limit / target-side correspondence
  theorem.} No PDE identification, no metric / spectral / Hecke
  / cohomology correspondence.
\item \textbf{Selection mechanism in the strong sense.}
  ``Cascade closure as selection mechanism'' is overstatement.
  This paper provides a top-rung $(O1)$ witness for a chosen
  $\xi^{*}$ on the $H_4$ branch at finite $N$, no more.
\end{enumerate}
\end{remark}

\subsection{Inputs and notation}

We use the following inputs without re-deriving:
\begin{itemize}
\item \cite{RefinementSpaces} (paper P3): the finite-level
  vertex-sector Hilbert space
  $X_n^{0, \bullet} = \ell^2(V(G_n^\bullet); F^0)$
  (\texttt{def:Xn-0}, \texttt{prop:Xn-Hilbert}); the
  Schl\"{a}fli-style refinement step (\texttt{def:refinement});
  the base graph $G_0^{H_4}$ on $V_{600}$
  (\texttt{def:base-graphs}); the sector fibres $F^0 = V_{H_4}
  \oplus V_{D_4} \oplus V_{16} \oplus V_8$ of $\R$-dimension
  $32$ (\texttt{not:fibres}); the branchwise Coxeter action
  $\rho_{F^0}^\bullet$ (\texttt{def:coxeter}); Coxeter unitarity
  (\texttt{lem:coxeter-unitary}); the vertex bonding map $p^0$
  (\texttt{def:p0}, \texttt{thm:bonding}).
\item \cite{ClosureDynamics} (paper P4): the discrete coboundary
  $d_n$ which projects $F^0$ to $F^1 = \mathrm{Im}(\OO)$ via
  $\pi_{F^0 \to F^1}$ (\texttt{def:coboundary}); the
  \emph{vertex} (top-left) block $A_n^{\mathrm{vertex}} :=
  d_n^{*} d_n$ on $X_n^{0, \bullet}$ (\texttt{def:An}); the
  finite-dim ODE / matrix-exponential argument for flow
  existence (\texttt{thm:flow-exists}).
\item \cite{RefinementCompat}: the abstract scalar pure-midpoint
  refinement model (\texttt{def:abstract-model}); the scaled
  curvature operator $\widetilde A_n := 2^n A_n^{\mathrm{vertex}}$
  in that abstract scalar model (\texttt{def:Atilde}); the
  open lift hypotheses $(L1), (L3)$
  (\texttt{sec:scope-and-gap}).
\item \cite{CascadeMechanism}: the cascade-closure consumer API
  $(O0)$--$(O3)$ (\S 3.4 lines 594--627); the explicit
  ``selection is the successor's task'' framing (\S 3.4 lines
  637--639).
\item \cite{AlgebraicSubstrate} (paper P2): the $V_{600}$
  vertex set (\texttt{def:V600}); the icosian ring
  (\texttt{def:icosian-ring}); $V_{24} \subset V_{600}$
  (\texttt{prop:24-in-600}); the explicit $V_{600}$ graph
  Laplacian spectrum (\texttt{tab:P2-spectrum}), used only in
  the $N = 0$ worked numerics of \S\ref{sec:numerics}.
\end{itemize}

Classical references invoked: \cite{Humphreys} for $H_4$
reflection-group irreducibility; \cite{Chung} for the scalar
graph Laplacian on multigraphs; \cite{ReedSimonI},
\cite{HornJohnson} for finite-dim spectral / matrix-exponential
arguments; \cite{CoxeterRegularPolytopes} for the $600$-cell edge
graph.

% =====================================================================
\section{The auxiliary block graph Laplacian $L_N^{\mathrm{block}}$}
\label{sec:block-laplacian}
% =====================================================================

\subsection{Why we need a new operator}

The cascade-closure consumer API of
\cite[\S 3.4]{CascadeMechanism} requires the supplied $R_N$ and
flow $\Psi_N^{t}$ to be the successor paper's own choice; in
particular, $R_N$ is not required to be P4's closure functional
$F_n = \alpha R_n + \beta E_n - \gamma Q_n$
(\cite[Definition~\texttt{def:Fn}]{ClosureDynamics}). We exploit
this freedom: P4's vertex / top-left block
$A_n^{\mathrm{vertex}} := d_n^{*} d_n$
(\cite[Definition~\texttt{def:An}]{ClosureDynamics}) acts
non-trivially only on the $\mathrm{Im}(\OO)$ direction within the
$V_8$ fibre summand, because P4's coboundary
$d_n \colon X_n^0 \to X_n^1$ first projects $F^0$ onto
$F^1 = \mathrm{Im}(\OO)$ via the canonical projection
$\pi_{F^0 \to F^1}$
(\cite[Definition~\texttt{def:coboundary}]{ClosureDynamics}). Hence
\[
  \dim_{\R} \ker A_n^{\mathrm{vertex}}
  \;=\; 25 \,|V(G_n^\bullet)| + 7
\]
on a connected graph $G_n^\bullet$ (the $25$ counts $\dim
V_{H_4} + \dim V_{D_4} + \dim V_{16} = 4 + 4 + 16$ plus the
real-octonion direction $\dim \R \cdot 1 = 1$, all of which
$d_n$ ignores; the additional $+7$ counts constant
$\mathrm{Im}(\OO)$-functions). That kernel is too large for a
clean $(O1)$ witness.

We need an operator whose kernel is exactly the $F^0$-valued
constants on a connected refined graph (dimension $|F^0| = 32$);
the simple block-diagonal extension of the scalar graph
Laplacian to $F^0$-valued functions does the job and is what
we introduce below.

\subsection{Definition}

\begin{definition}[Labelled-multiset scalar graph Laplacian]
\label{def:scalar-laplacian}
Let $G_N^{H_4} = (V(G_N^{H_4}), E(G_N^{H_4}))$ be the level-$N$
Schl\"{a}fli-style refined graph of
\cite[Definition~\texttt{def:refinement}]{RefinementSpaces},
with base $G_0^{H_4}$ on $V_{600}$
(\cite[Definition~\texttt{def:base-graphs}]{RefinementSpaces}).
By \cite[Definition~\texttt{def:refinement}]{RefinementSpaces},
$E(G_N^{H_4})$ is a \emph{labelled-edge multiset}: each generated
edge carries a class tag $\in \{(\mathrm{a}), (\mathrm{b}),
(\mathrm{c})\}$ and a parent-face tag where applicable, so two
edges with the same unordered endpoint pair but different tags
are distinct elements of $E(G_N^{H_4})$.

The \emph{labelled-multiset combinatorial scalar graph Laplacian}
$L^{\mathrm{cl}}_N$ on $V(G_N^{H_4})$, acting on $\R$-valued
functions, is
\[
  (L^{\mathrm{cl}}_N f)(v)
  \;:=\; \sum_{e \in E(G_N^{H_4}),\; v \in e}
         \bigl(f(v) - f(\op_e(v))\bigr),
\]
where for each $e = \{v, \op_e(v)\}$ in the multiset, $\op_e(v)$
denotes the other endpoint of the labelled edge $e$. Each labelled
edge $e \in E(G_N^{H_4})$ contributes once with multiplicity $1$;
parallel labelled edges with coinciding endpoints are summed
separately.

Equivalently, $L^{\mathrm{cl}}_N = D_N - W_N$ where the diagonal
$D_N(v) := |\{e \in E(G_N^{H_4}) : v \in e\}| =
\deg_{G_N^{H_4}}(v)$ uses the labelled-multiset degree count, and
$W_N(v, u) := |\{e \in E(G_N^{H_4}) : e = \{v, u\}\}|$ is the
number of labelled edges joining $v$ to $u$ (an integer
$\geq 0$). This is the standard Chung convention
(\cite[Ch.~1]{Chung}) extended to multigraphs.

For $N = 0$ all base edges of $G_0^{H_4}$ on $V_{600}$ have
multiplicity $1$ (the $600$-cell edge graph is simple,
\cite[\S 22.7]{CoxeterRegularPolytopes}), so $W_0(v, u) \in
\{0, 1\}$ and $L^{\mathrm{cl}}_0$ reduces to the standard simple
graph Laplacian on the $600$-cell edge graph.
\end{definition}

\begin{definition}[Block graph Laplacian]
\label{def:block-laplacian}
The \emph{block graph Laplacian on the $H_4$ branch at level $N$}
is the bounded $\R$-linear operator $L_N^{\mathrm{block}}$ on the
finite-dimensional Hilbert space
$X_N^{0, H_4} = \ell^2(V(G_N^{H_4}); F^0)$
(\cite[Definition~\texttt{def:Xn-0}]{RefinementSpaces}) defined
fibrewise by
\[
  L_N^{\mathrm{block}}(\phi)(v)
  \;:=\; \bigl(L^{\mathrm{cl}}_N \phi^{(H_4)}\bigr)(v)
        \oplus \bigl(L^{\mathrm{cl}}_N \phi^{(D_4)}\bigr)(v)
        \oplus \bigl(L^{\mathrm{cl}}_N \phi^{(16)}\bigr)(v)
        \oplus \bigl(L^{\mathrm{cl}}_N \phi^{(8)}\bigr)(v),
\]
where $\phi(v) = \phi^{(H_4)}(v) \oplus \phi^{(D_4)}(v) \oplus
\phi^{(16)}(v) \oplus \phi^{(8)}(v)$ in the orthogonal
decomposition $F^0 = V_{H_4} \oplus V_{D_4} \oplus V_{16} \oplus
V_8$ of \cite[Notation~\texttt{not:fibres}]{RefinementSpaces},
and $L^{\mathrm{cl}}_N$ acts componentwise on each summand
independently via Definition~\ref{def:scalar-laplacian}.
\end{definition}

\subsection{Relation to upstream operators}

\begin{remark}[$L_N^{\mathrm{block}}$ is not an upstream operator]
\label{rmk:no-identification}
We do \emph{not} claim that $L_N^{\mathrm{block}}$ equals or is a
restriction of any operator from \cite{ClosureDynamics} or
\cite{RefinementCompat}. Specifically:
\begin{itemize}
\item P4's vertex / top-left block
  $A_n^{\mathrm{vertex}} = d_n^{*} d_n$
  (\cite[Definition~\texttt{def:An}]{ClosureDynamics}) factors
  through the $\mathrm{Im}(\OO)$ projection of P4's coboundary
  (\cite[Definition~\texttt{def:coboundary}]{ClosureDynamics}):
  informally,
  $A_n^{\mathrm{vertex}} = \iota_{F^1 \to F^0} \circ
  L^{\mathrm{cl}}_N \circ \pi_{F^0 \to F^1}$. Compared with
  $L_N^{\mathrm{block}}$, this replaces three of the four fibre
  blocks ($V_{H_4}$, $V_{D_4}$, $V_{16}$, plus the real-octonion
  direction in $V_8$) with zero. Hence $L_N^{\mathrm{block}}$ is
  a strictly larger operator on $X_N^{0, H_4}$.
\item \cite{RefinementCompat}'s scaled curvature operator
  $\widetilde A_n := 2^n A_n^{\mathrm{vertex}}$
  (\cite[Definition~\texttt{def:Atilde}]{RefinementCompat}) is
  the \emph{scaled scalar abstract-model analogue} of P4's
  vertex Laplacian: \cite{RefinementCompat} works in the
  abstract scalar pure-midpoint reduction
  (\cite[Definition~\texttt{def:abstract-model}]{RefinementCompat}),
  where the underlying state space is scalar-valued (not
  $F^0$-valued) on a pure-midpoint refinement of $V_{600}$. The
  full-fibre $L_N^{\mathrm{block}}$ on $X_N^{0, H_4}$ used here
  is in a different scope from both
  \cite{RefinementCompat} (which has no fibre $F^0$) and
  \cite{ClosureDynamics} (which restricts to the
  $\mathrm{Im}(\OO)$ direction).
\end{itemize}
$L_N^{\mathrm{block}}$ is introduced in this paper for the
$(O1)$ witness. The simple block-diagonal construction is one
natural choice; we do not claim minimality.
\end{remark}

% =====================================================================
\section{Self-adjointness, positive semi-definiteness, and bound}
\label{sec:psd}
% =====================================================================

\begin{proposition}[Clause (1): operator structure of
$L_N^{\mathrm{block}}$]
\label{prop:psd}
For every $N \geq 0$, $L_N^{\mathrm{block}}$ is bounded,
self-adjoint, and positive semi-definite on $X_N^{0, H_4}$, with
operator-norm bound
\[
  \|L_N^{\mathrm{block}}\|_{\op}
  \;\leq\; 2 \, \Delta(G_N^{H_4})
\]
where $\Delta(G_N^{H_4}) := \max_{v \in V(G_N^{H_4})}
\deg_{G_N^{H_4}}(v)$ is the maximum labelled-multiset degree.
\end{proposition}

\begin{proof}
Self-adjointness and positive semi-definiteness of the scalar
Laplacian $L^{\mathrm{cl}}_N$ follow from the standard identity
\[
  \langle L^{\mathrm{cl}}_N f, g \rangle_{\ell^2(V(G_N^{H_4}); \R)}
  \;=\; \sum_{e \in E(G_N^{H_4})}
        (f(h(e)) - f(t(e))) (g(h(e)) - g(t(e)))
\]
(for any orientation $e \mapsto (t(e), h(e))$ of each labelled
edge; the sum is symmetric in the unoriented edge), which gives
$\langle L^{\mathrm{cl}}_N f, g\rangle = \langle f,
L^{\mathrm{cl}}_N g\rangle$ and
$\langle L^{\mathrm{cl}}_N f, f\rangle = \sum_e (f(h(e)) -
f(t(e)))^{2} \geq 0$ (\cite[Ch.~1]{Chung} for the multigraph
extension).

The block-diagonal construction
$L_N^{\mathrm{block}} = L^{\mathrm{cl}}_N \otimes \id_{F^0}$
(after identifying $X_N^{0, H_4} = \ell^2(V(G_N^{H_4}); \R)
\otimes_{\R} F^0$ orthogonally) preserves both self-adjointness
and positive semi-definiteness.

For the bound, the standard scalar estimate
$\|L^{\mathrm{cl}}_N\|_{\op} \leq 2 \, \Delta(G_N^{H_4})$
(\cite[Ch.~1]{Chung}, multigraph case) extends to the block
operator with the same constant via the tensor-product
operator-norm identity: under the orthogonal identification
$X_N^{0, H_4} \cong \ell^2(V(G_N^{H_4}); \R) \otimes_{\R} F^{0}$
of Definition~\ref{def:block-laplacian}, $L_N^{\mathrm{block}} =
L^{\mathrm{cl}}_N \otimes \id_{F^{0}}$, so
$\|L_N^{\mathrm{block}}\|_{\op} = \|L^{\mathrm{cl}}_N\|_{\op}
\cdot \|\id_{F^{0}}\|_{\op} = \|L^{\mathrm{cl}}_N\|_{\op}$
(\cite[Ch.~4.2]{HornJohnson}).
\end{proof}

% =====================================================================
\section{Connectivity and the kernel theorem}
\label{sec:kernel}
% =====================================================================

\begin{lemma}[Connectivity of the refined $H_4$ graph]
\label{lem:connected}
For every $N \geq 0$, the (labelled-multiset) graph $G_N^{H_4}$
is connected.
\end{lemma}

\begin{proof}
Induction on $N$.

\emph{Base case $N = 0$.} $G_0^{H_4}$ has vertex set $V_{600}$
and edges joining $600$-cell vertices at the standard $600$-cell
edge distance
(\cite[Definition~\texttt{def:base-graphs}]{RefinementSpaces},
\cite[\S 22.7]{CoxeterRegularPolytopes}). The $600$-cell edge
graph is connected; this is classical and equivalent to
$W(H_4)$-orbit transitivity on $V_{600}$ via edge-adjacent unit
icosians.

\emph{Inductive step.} Assume $G_n^{H_4}$ is connected. By
\cite[Definition~\texttt{def:refinement}]{RefinementSpaces},
$V(G_{n+1}^{H_4}) = V(G_n^{H_4}) \sqcup M_n^{H_4} \sqcup
C_n^{H_4}$ (vertices, midpoints, centroids), and the edge set
$E(G_{n+1}^{H_4})$ consists of three classes:
\begin{itemize}[leftmargin=2em]
\item Class (a): for each $e = \{u, v\} \in E_n^{H_4}$, two
  subdividing edges $\{u, m_e\}$ and $\{m_e, v\}$ (using the
  midpoint $m_e$).
\item Class (b): for each face $f \in F_n^{H_4}$ with edge
  sequence $(e_1, \ldots, e_k)$, midpoint-centroid edges
  $\{m_{e_i}, c_f\}$.
\item Class (c): midpoint-midpoint diagonal edges
  $\{m_{e_i}, m_{e_{i+1}}\}$ closing each level-$(n+1)$ triangle
  face.
\end{itemize}
Class-(a) edges keep every level-$(n+1)$ midpoint $m_e$ joined to
the existing $V(G_n^{H_4})$ component (each $m_e$ is adjacent to
both endpoints of its parent edge $e$). Class-(b) edges keep
every level-$(n+1)$ centroid $c_f$ joined to the existing
component (each $c_f$ is adjacent to every midpoint of every
edge in $\partial^\bullet_n(f)$, all of which lie in
$V(G_{n+1}^{H_4})$ via Class (a)). Hence every vertex in
$V(G_{n+1}^{H_4})$ is in the same connected component as
$V(G_n^{H_4})$, which is connected by the inductive hypothesis.
Class-(c) edges are not needed for the inductive step.
\end{proof}

\begin{theorem}[Clause (2): kernel structure of
$L_N^{\mathrm{block}}$]
\label{thm:kernel}
For every $N \geq 0$, the constant-function embedding
\[
  \iota_N \colon F^0 \to \ker L_N^{\mathrm{block}},
  \qquad \iota_N(\xi)(v) := \xi \quad \text{for all }
  v \in V(G_N^{H_4}),
\]
is an $\R$-linear isomorphism. In particular,
$\dim_{\R} \ker L_N^{\mathrm{block}} = \dim_{\R} F^0 = 32$.
\end{theorem}

\begin{proof}
The kernel of the scalar Laplacian on a connected (multi)graph
is one-dimensional, spanned by the constant function
(\cite[Ch.~1]{Chung}, multigraph case). By Lemma~\ref{lem:connected},
$G_N^{H_4}$ is connected, so $\ker L^{\mathrm{cl}}_N = \R \cdot
\mathbf{1}_{V(G_N^{H_4})}$.

By the tensor identification of Definition~\ref{def:block-laplacian},
$L_N^{\mathrm{block}} = L^{\mathrm{cl}}_N \otimes \id_{F^{0}}$ on
$X_N^{0, H_4} \cong \ell^2(V(G_N^{H_4}); \R) \otimes_{\R} F^{0}$,
so
\[
  \ker L_N^{\mathrm{block}}
  \;=\; \bigl(\ker L^{\mathrm{cl}}_N\bigr) \otimes_{\R} F^{0}
  \;=\; \bigl(\R \cdot \mathbf{1}_{V(G_N^{H_4})}\bigr) \otimes_{\R} F^{0}
  \;\cong\; F^{0},
\]
where the last isomorphism sends $\mathbf{1} \otimes \xi$ to
$\xi \in F^{0}$. This isomorphism is exactly $\iota_N$
(constant-function embedding $\xi \mapsto (v \mapsto \xi)$).

$\R$-linearity of $\iota_N$ is by construction; injectivity is
immediate; surjectivity onto $\ker L_N^{\mathrm{block}}$ is the
above identification.
\end{proof}

% =====================================================================
\section{Convergence theorem}
\label{sec:convergence}
% =====================================================================

\begin{notation}[Spectral gap]
\label{not:lambda-plus}
For a positive semi-definite self-adjoint operator $T$ on a
finite-dimensional Hilbert space, write
\[
  \lambda_{+}(T)
  \;:=\; \min\bigl\{\lambda \in \spec(T) : \lambda > 0\bigr\}
\]
for the smallest positive eigenvalue (the spectral gap). When
$\dim \ker T > 1$, $\lambda_{+}(T)$ is strictly larger than the
``second eigenvalue'' $\lambda_2(T) = 0$, and we use $\lambda_{+}$
throughout to avoid confusion.
\end{notation}

\begin{theorem}[Clause (3): convergence of the gradient flow]
\label{thm:convergence}
For every $N \geq 0$ and every $\Phi(0) \in X_N^{0, H_4}$, the
operator semigroup $e^{-t L_N^{\mathrm{block}}}$ satisfies
\[
  \lim_{t \to \infty} e^{-t L_N^{\mathrm{block}}} \Phi(0)
  \;=\; P_{\ker L_N^{\mathrm{block}}}\,\Phi(0)
  \quad \text{in } X_N^{0, H_4}\text{-norm,}
\]
with exponential bound, uniform over initial data $\Phi(0)$ at
fixed $N$ (no claim about uniformity in $N$),
\[
  \bigl\|\, e^{-t L_N^{\mathrm{block}}} \Phi(0)
        - P_{\ker L_N^{\mathrm{block}}}\Phi(0) \,\bigr\|_{X_N^{0, H_4}}
  \;\leq\;
  e^{-t \, \lambda_{+}(L_N^{\mathrm{block}})}
  \,\bigl\|\Phi(0) - P_{\ker L_N^{\mathrm{block}}}\Phi(0)\bigr\|_{X_N^{0, H_4}}
\]
for every $t \geq 0$, where $\lambda_{+}(L_N^{\mathrm{block}}) =
\lambda_{+}(L^{\mathrm{cl}}_N) > 0$ (positivity is by
Lemma~\ref{lem:connected}: the connected, nontrivial multigraph
$G_N^{H_4}$ has a strictly positive Fiedler value; the block
extension repeats each scalar eigenvalue with multiplicity $32$,
leaving the smallest positive eigenvalue unchanged).
\end{theorem}

\begin{proof}
By Proposition~\ref{prop:psd}, $L_N^{\mathrm{block}}$ is bounded
self-adjoint and positive semi-definite on the finite-dimensional
Hilbert space $X_N^{0, H_4}$. By the spectral theorem
(\cite[Ch.~VIII.5]{ReedSimonI},
\cite[Ch.~6.2]{HornJohnson}), there is an orthogonal decomposition
\[
  X_N^{0, H_4} = \bigoplus_{\lambda \in \spec(L_N^{\mathrm{block}})}
                  E_\lambda(L_N^{\mathrm{block}}),
\]
where $E_0 = \ker L_N^{\mathrm{block}}$ and every other
$E_\lambda$ has $\lambda \geq \lambda_{+}(L_N^{\mathrm{block}})$.
For $\Phi(0) = \sum_\lambda \Phi_\lambda(0)$ with
$\Phi_\lambda(0) \in E_\lambda$,
\[
  e^{-t L_N^{\mathrm{block}}} \Phi(0)
  \;=\; \sum_\lambda e^{-t\lambda} \Phi_\lambda(0)
  \;=\; \Phi_0(0) + \sum_{\lambda > 0} e^{-t\lambda} \Phi_\lambda(0).
\]
The first term is $P_{\ker} \Phi(0)$, constant in $t$. The
second term is bounded by
\[
  \Bigl\| \sum_{\lambda > 0} e^{-t\lambda} \Phi_\lambda(0)
  \Bigr\|^{2}
  = \sum_{\lambda > 0} e^{-2t\lambda} \|\Phi_\lambda(0)\|^{2}
  \leq e^{-2t\lambda_{+}}
       \sum_{\lambda > 0} \|\Phi_\lambda(0)\|^{2}
  = e^{-2t\lambda_{+}}
    \|\Phi(0) - P_{\ker}\Phi(0)\|^{2},
\]
which gives the stated exponential bound on taking square roots.

The eigenvalue identity $\lambda_{+}(L_N^{\mathrm{block}}) =
\lambda_{+}(L^{\mathrm{cl}}_N)$ follows from
$\spec(L_N^{\mathrm{block}}) = \spec(L^{\mathrm{cl}}_N)$
(set-equality; multiplicities multiply by $\dim F^0 = 32$ but the
set of distinct eigenvalues is unchanged), which is a standard
property of block-diagonal operators
(\cite[Ch.~6.4]{HornJohnson}).
\end{proof}

% =====================================================================
\section{The $W(H_4)$-fixed subspace count}
\label{sec:fixed}
% =====================================================================

\begin{lemma}[Standard $H_4$-reflection representation has trivial
fixed subspace]
\label{lem:H4-fixed}
The standard reflection representation of the finite irreducible
Coxeter group $W(H_4)$ on $V_{H_4} \cong \HH \cong \R^{4}$
satisfies $V_{H_4}^{W(H_4)} = \{0\}$.
\end{lemma}

\begin{proof}
$W(H_4)$ is a finite \emph{irreducible} Coxeter group; the
standard reflection representation of any finite irreducible
Coxeter group is irreducible (\cite[Ch.~5]{Humphreys}; see also
\cite[Ch.~2]{Humphreys} for $H_4$ specifically). The reflection
representation of a finite irreducible Coxeter group is faithful
and is not the trivial representation
(\cite[Ch.~5]{Humphreys}). The fixed subspace
$V_{H_4}^{W(H_4)} \subset V_{H_4}$ is a $W(H_4)$-invariant
subspace; by irreducibility it is either $\{0\}$ or all of
$V_{H_4}$. The latter would force every reflection to act as
the identity, contradicting faithfulness (or non-triviality).
Hence $V_{H_4}^{W(H_4)} = \{0\}$.
\end{proof}

\begin{theorem}[Clause (4): the $W(H_4)$-fixed subspace of $F^0$]
\label{thm:fixed-dim}
Under the branchwise Coxeter representation $\rho_{F^0}^{H_4}$ of
\cite[Definition~\texttt{def:coxeter}]{RefinementSpaces}
(standard $H_4$-reflection representation on $V_{H_4}$, trivial
on $V_{D_4}, V_{16}, V_8$),
\[
  (F^0)^{W(H_4)}
  \;=\; V_{H_4}^{W(H_4)} \oplus V_{D_4} \oplus V_{16} \oplus V_8
  \;=\; \{0\} \oplus V_{D_4} \oplus V_{16} \oplus V_8,
\]
so $\dim_{\R} (F^0)^{W(H_4)} = 0 + 4 + 16 + 8 = 28$.
\end{theorem}

\begin{proof}
The branchwise representation
$\rho_{F^0}^{H_4} = \rho_{V_{H_4}}^{H_4} \oplus
\rho_{V_{D_4}}^{H_4} \oplus \rho_{V_{16}}^{H_4} \oplus
\rho_{V_8}^{H_4}$ acts on each summand independently
(\cite[Definition~\texttt{def:coxeter}]{RefinementSpaces} fixes
the standard reflection representation on $V_{H_4}$ and the
trivial representation on the other three summands). Hence
\[
  (F^0)^{W(H_4)}
  = V_{H_4}^{W(H_4)} \oplus V_{D_4}^{W(H_4)} \oplus
    V_{16}^{W(H_4)} \oplus V_{8}^{W(H_4)}.
\]
By Lemma~\ref{lem:H4-fixed}, $V_{H_4}^{W(H_4)} = \{0\}$.
By the trivial representation on the other summands,
$V_{D_4}^{W(H_4)} = V_{D_4}$ (dim $4$),
$V_{16}^{W(H_4)} = V_{16}$ (dim $16$),
$V_{8}^{W(H_4)} = V_{8}$ (dim $8$). Summing dimensions,
$\dim_{\R} (F^0)^{W(H_4)} = 0 + 4 + 16 + 8 = 28$.
\end{proof}

% =====================================================================
\section{$(O0)$ verification and $(O1)$ witnessed discharge}
\label{sec:o1-discharge}
% =====================================================================

\subsection{$(O0)$ standing analytic preconditions}

\begin{proposition}[Top-rung analytic preconditions and recorded
inter-rung bonding map]
\label{prop:o0}
The top-rung analytic components of the cascade-mechanism standing
preconditions $(O0)$ of
\cite[\S 3.4 lines 595--608]{CascadeMechanism} hold for the
five-tuple supplied in Theorem~\ref{thm:main} below; the
inter-rung bonding map needed by the rung-sequence formulation of
$(O0)$ is recorded as available from
\cite[Definition~\texttt{def:p0}, Theorem~\texttt{thm:bonding}]
{RefinementSpaces} but is not used here. Specifically:
\begin{enumerate}[label=\textup{(O0.\arabic*)}, leftmargin=3em]
\item \emph{Finite-dimensional Hilbert structure of $\Phi_N =
  X_N^{0, H_4}$.} By
  \cite[Proposition~\texttt{prop:Xn-Hilbert}]{RefinementSpaces},
  $X_N^{0, H_4}$ is a finite-dimensional real Hilbert space of
  dimension $|V(G_N^{H_4})| \cdot 32$.
\item \emph{$R_N \in C^{1}(\Phi_N)$.} The quadratic form
  $R_N(\phi) = \tfrac12 \langle \phi, L_N^{\mathrm{block}} \phi
  \rangle$ on a finite-dimensional Hilbert space is real-analytic
  (in particular $C^{1}$); its gradient is $\nabla R_N(\phi) =
  L_N^{\mathrm{block}} \phi$ by Proposition~\ref{prop:psd}
  (self-adjointness).
\item \emph{Well-posed gradient flow $\Psi_N^{t} = e^{-t
  L_N^{\mathrm{block}}}$.} By the same finite-dimensional ODE /
  matrix-exponential argument as
  \cite[Theorem~\texttt{thm:flow-exists}]{ClosureDynamics},
  applied here to the bounded self-adjoint $L_N^{\mathrm{block}}$
  on the finite-dimensional Hilbert space $X_N^{0, H_4}$, the
  initial-value problem $\dot\Phi_N(t) = -L_N^{\mathrm{block}}
  \Phi_N(t)$, $\Phi_N(0) = \Phi(0)$ has a unique global solution
  $\Phi_N(t) = e^{-t L_N^{\mathrm{block}}} \Phi(0)$.
\item \emph{$C^{1}$ pushforward bonding maps (recorded only).} The
  $(O1)$ witness below uses only top-rung residual-zero and
  convergence data; no inter-rung pushforward is invoked. For
  completeness, when the cascade-mechanism API invokes an
  inter-rung pushforward $\pi^{\sharp}_{k+1, k} \colon \Phi_{k+1}
  \to \Phi_k$ for downstream propagation
  (\cite[\S 3.4 lines 600--605]{CascadeMechanism}),
  \cite[Definition~\texttt{def:p0}, Theorem~\texttt{thm:bonding}]
  {RefinementSpaces} supplies a bounded $\R$-linear (hence $C^{1}$,
  in fact $C^{\omega}$) vertex restriction map
  $p_{k+1, k}^{0} \colon X_{k+1}^{0, H_4} \to X_k^{0, H_4}$
  between adjacent levels. This paper does not use it for the
  $(O1)$ discharge but records its availability for downstream
  consumers wishing to instantiate the full rung-sequence
  formulation of $(O0)$.
\end{enumerate}
\end{proposition}

\begin{proof}
Each item follows directly from the cited result, as stated.
\end{proof}

\subsection{$(O1)$ witnessed discharge: the named theorem}

\begin{theorem}[Witnessed $(O1)$ discharge on the $H_4$ branch at
finite top-rung level $N$, parametrised by $\xi^{*}$]
\label{thm:main}
Let $N \geq 0$ and let $L_N^{\mathrm{block}}$ be the block graph
Laplacian of Definition~\ref{def:block-laplacian}. For each
fixed
\[
  \xi^{*} \in (F^0)^{W(H_4)} \setminus \{0\},
\]
where $(F^0)^{W(H_4)}$ is the $W(H_4)$-fixed subspace of $F^0$ of
explicit real dimension $28$ (Theorem~\ref{thm:fixed-dim}), set:
\begin{itemize}
\item $\Phi_N := X_N^{0, H_4}$
  (\cite[Definition~\texttt{def:Xn-0}]{RefinementSpaces});
\item $R_N(\phi) := \tfrac{1}{2}
  \langle \phi, L_N^{\mathrm{block}} \phi \rangle_{X_N^{0, H_4}}$,
  the quadratic-form residual associated to $L_N^{\mathrm{block}}$;
\item $\Psi_N^{t} := e^{-t L_N^{\mathrm{block}}}$, the gradient
  flow of $R_N$;
\item $\phi^{\star}_{N} := \iota_N(\xi^{*})$, the constant
  $\xi^{*}$-valued function on $V(G_N^{H_4})$
  (Theorem~\ref{thm:kernel});
\item $\phi_N^{(0)} := \phi^{\star}_{N}$.
\end{itemize}
Then the five-tuple $(\Phi_N, R_N, \Psi_N^{t}, \phi^{\star}_{N},
\phi_N^{(0)})$ satisfies the cascade-mechanism $(O1)$ obligation
of \cite[\S 3.4 lines 609--617]{CascadeMechanism} for the chosen
$\xi^{*}$ and the chosen top-rung initial datum
$\phi_N^{(0)} = \phi^{\star}_{N}$:
\begin{enumerate}[label=\textup{(O1.\arabic*)}, leftmargin=3em]
\item \emph{Top-rung residual-zero state.} $\phi^{\star}_{N} \in
  \Phi_N$ satisfies $R_N(\phi^{\star}_{N}) = 0$.
\item \emph{Convergence from a chosen initial datum.} For
  $\phi_N^{(0)} = \phi^{\star}_{N}$, $\Psi_N^{t} \phi_N^{(0)} =
  \phi^{\star}_{N}$ for every $t \geq 0$ (a fixed point of the
  flow). More generally, for any $\widetilde\phi_N^{(0)} \in
  \Phi_N$ with $P_{\ker L_N^{\mathrm{block}}} \widetilde
  \phi_N^{(0)} = \phi^{\star}_{N}$,
  $\lim_{t \to \infty} \Psi_N^{t} \widetilde\phi_N^{(0)} =
  \phi^{\star}_{N}$ in $X_N^{0, H_4}$-norm with uniform
  exponential bound governed by
  $\lambda_{+}(L_N^{\mathrm{block}})$
  (Theorem~\ref{thm:convergence}).
\end{enumerate}
The standing analytic preconditions $(O0)$ are verified in
Proposition~\ref{prop:o0}.
\end{theorem}

\begin{proof}
\emph{(O1.1).} By Theorem~\ref{thm:kernel}, $\iota_N(\xi^{*}) \in
\ker L_N^{\mathrm{block}}$, so
$L_N^{\mathrm{block}} \phi^{\star}_{N} = 0$, hence
$R_N(\phi^{\star}_{N}) =
\tfrac12 \langle \phi^{\star}_{N}, L_N^{\mathrm{block}}
\phi^{\star}_{N} \rangle = 0$.

\emph{(O1.2).} For $\phi_N^{(0)} = \phi^{\star}_{N}$,
$\Psi_N^{t} \phi^{\star}_{N} = e^{-t L_N^{\mathrm{block}}}
\phi^{\star}_{N} = \phi^{\star}_{N}$ for all $t \geq 0$ because
$\phi^{\star}_{N} \in \ker L_N^{\mathrm{block}}$ implies
$L_N^{\mathrm{block}} \phi^{\star}_{N} = 0$, hence
$e^{-t L_N^{\mathrm{block}}} \phi^{\star}_{N} =
\sum_{k \geq 0} \frac{(-t)^{k}}{k!}
(L_N^{\mathrm{block}})^{k} \phi^{\star}_{N} = \phi^{\star}_{N}$
(only the $k = 0$ term survives).

For the more general statement: any $\widetilde\phi_N^{(0)}$ with
$P_{\ker} \widetilde\phi_N^{(0)} = \phi^{\star}_{N}$ decomposes
orthogonally as
$\widetilde\phi_N^{(0)} = \phi^{\star}_{N} + \widetilde\phi^{\perp}$
with $\widetilde\phi^{\perp} \in (\ker L_N^{\mathrm{block}})^{\perp}$.
By Theorem~\ref{thm:convergence},
$\Psi_N^{t} \widetilde\phi^{\perp} \to 0$ in norm with rate
$\lambda_{+}$, while $\Psi_N^{t} \phi^{\star}_{N} =
\phi^{\star}_{N}$ as above. Hence
$\Psi_N^{t} \widetilde\phi_N^{(0)} \to \phi^{\star}_{N}$ in norm.
\end{proof}

\begin{remark}[On the $\xi^{*}$ parameter]
\label{rmk:xi-freedom}
For each $\xi^{*} \in (F^0)^{W(H_4)} \setminus \{0\}$,
Theorem~\ref{thm:main} produces a witnessed $(O1)$-tuple. The
freedom in $\xi^{*}$ is the $(O1)$-witness parameter; the
\cite[\S 3.4 lines 637--639]{CascadeMechanism} statement
``selection of $\phi^{\star}_{N}$ itself is not provided by
cascade closure and is the successor paper's own task'' is
satisfied by our supplying \emph{a} $\phi^{\star}_{N}$, not the
unique one. Reducing the $28$-dimensional $\xi^{*}$-freedom to a
single ray requires the named open hypothesis
\textbf{(H-uniqueness)} (\S\ref{sec:open-hyp}).

The chosen $\xi^{*}$ is $W(H_4)$-equivariant by construction
($\xi^{*} \in (F^0)^{W(H_4)}$), so $\iota_N(\xi^{*})$ is a
$W(H_4)$-equivariant function on $V(G_N^{H_4})$.
$W(H_4)$-non-equivariant initial data
$\widetilde\phi_N^{(0)} \in X_N^{0, H_4}$ still converge by
Theorem~\ref{thm:convergence}, but to the $32$-dimensional full kernel
projection $P_{\ker L_N^{\mathrm{block}}} \widetilde\phi_N^{(0)}$,
which need not coincide with any constant
$\iota_N(\xi^{*})$ for $\xi^{*}$ in the $28$-dimensional
$W(H_4)$-fixed subspace.
\end{remark}

% =====================================================================
\section{Worked numerics for $N = 0$}
\label{sec:numerics}
% =====================================================================

We exhibit the explicit numerical content of
Theorem~\ref{thm:main} at the base level $N = 0$, where the
graph $G_0^{H_4}$ on $V_{600}$ has a fully understood spectrum
imported from \cite[Table~\texttt{tab:P2-spectrum}]{AlgebraicSubstrate}.

\begin{itemize}[leftmargin=2em]
\item Vertex count: $|V(G_0^{H_4})| = |V_{600}| = 120$
  (\cite[Definition~\texttt{def:V600}]{AlgebraicSubstrate}).
\item Edge count: $|E(G_0^{H_4})| = 720$
  (the $600$-cell edge graph has $120$ vertices each of degree
  $12$, giving $\tfrac{120 \cdot 12}{2} = 720$ edges;
  \cite[\S 22.7]{CoxeterRegularPolytopes}).
\item Maximum labelled-multiset degree (here equal to ordinary
  degree since the base graph is simple):
  $\Delta(G_0^{H_4}) = 12$ (uniform).
\item Hilbert-space dimension:
  $\dim_{\R} X_0^{0, H_4} = 120 \cdot 32 = 3840$
  (\cite[Proposition~\texttt{prop:Xn-Hilbert}]{RefinementSpaces}).
\item Block-Laplacian operator-norm bound:
  $\|L_0^{\mathrm{block}}\|_{\op} \leq 2 \cdot 12 = 24$
  (Proposition~\ref{prop:psd}). The actual operator norm is
  $\|L_0^{\mathrm{block}}\|_{\op} = 9 + 3\sqrt{5} \approx 15.708$
  (largest eigenvalue from
  \cite[Table~\texttt{tab:P2-spectrum}]{AlgebraicSubstrate}).
\item Kernel dimension:
  $\dim_{\R} \ker L_0^{\mathrm{block}} = 32$
  (Theorem~\ref{thm:kernel}; the trivial irrep multiplicity is
  $1$ for the scalar Laplacian, so $1 \cdot 32 = 32$ in the
  block extension).
\item Spectral gap:
  $\lambda_{+}(L_0^{\mathrm{block}}) = \lambda_{+}(L^{\mathrm{cl}}_0)
  = 9 - 3\sqrt{5} \approx 2.2918$
  (smallest positive eigenvalue from
  \cite[Table~\texttt{tab:P2-spectrum}]{AlgebraicSubstrate}, the
  $\mathbf{2}$ irrep entry).
\item $W(H_4)$-fixed subspace of $F^0$:
  $\dim_{\R} (F^0)^{W(H_4)} = 28$ (Theorem~\ref{thm:fixed-dim}).
\item $\xi^{*}$-parametrised constant subspace inside
  $\ker L_0^{\mathrm{block}}$:
  $\dim_{\R} \iota_0\bigl((F^0)^{W(H_4)}\bigr) = 28$
  (image of injective $\iota_0$).
\end{itemize}

The convergence rate constant is therefore
\[
  \lambda_{+}(L_0^{\mathrm{block}}) = 9 - 3\sqrt{5},
\]
giving the explicit decay bound
$\bigl\|\Phi(t) - P_{\ker} \Phi(0)\bigr\|
\leq e^{-(9 - 3\sqrt{5}) t}
\bigl\|\Phi(0) - P_{\ker} \Phi(0)\bigr\|$
for any initial datum $\Phi(0) \in X_0^{0, H_4}$.

\begin{remark}[Sim verification (optional)]
\label{rmk:sim}
A one-line numerical check of the spectrum is straightforward:
build the $120 \times 120$ scalar Laplacian
$L^{\mathrm{cl}}_0 = D_0 - W_0$ from the $600$-cell adjacency
matrix and compute its eigenvalues. The result reproduces
\cite[Table~\texttt{tab:P2-spectrum}]{AlgebraicSubstrate}:
nine distinct eigenvalues $\{0, 9 - 3\sqrt{5}, 10 - 2\sqrt{5},
9, 12, 14, 10 + 2\sqrt{5}, 15, 9 + 3\sqrt{5}\}$ with
multiplicities $\{1, 4, 9, 16, 25, 36, 9, 16, 4\}$ summing to
$120$. The block-Laplacian spectrum is the same set of
eigenvalues, each with multiplicity multiplied by $32$.
\end{remark}

% =====================================================================
\section{What this paper does not do}
\label{sec:non-claims}
% =====================================================================

For convenience we collect the explicit non-claims of
Remark~\ref{rmk:non-claims} as bullet points, with cross-references
to where in this paper or in upstream sources each non-claim is
located:

\begin{itemize}[leftmargin=2em]
\item \textbf{Single-state selection.} See \S\ref{sec:open-hyp}
  \textbf{(H-uniqueness)} below; reducing the $28$-dimensional $\xi^{*}$
  freedom to a single ray is the next-paper task.
\item \textbf{$(O2)$ and $(O3)$.} Inherited from
  \cite{RefinementCompat}: $(O2)$ as stated remains false
  (\cite[Proposition~\texttt{prop:strict-fails}]{RefinementCompat}),
  and $(O3)$ holds only in the abstract scalar pure-midpoint
  model
  (\cite[Theorem~\texttt{thm:O3-discharge}]{RefinementCompat}).
\item \textbf{Inverse-limit $(O1)$.} Lifting Theorem~\ref{thm:main}
  to $N \to \infty$ requires the open hypotheses $(L1), (L3)$
  of \cite[\S\texttt{sec:scope-and-gap}]{RefinementCompat}.
\item \textbf{$D_4$ branch.} Theorem~\ref{thm:fixed-dim} and the
  surrounding construction use the $H_4$-branch action of
  \cite[Definition~\texttt{def:coxeter}]{RefinementSpaces}; the
  $D_4$ analogue would have a different fixed-subspace count and
  is a separate paper.
\item \textbf{Identification of $L_N^{\mathrm{block}}$ with
  upstream operators.} See Remark~\ref{rmk:no-identification}.
\item \textbf{$W(H_4)$-non-equivariant initial data convergent to
  $(F^0)^{W(H_4)}$.} See Remark~\ref{rmk:xi-freedom}.
\item \textbf{Continuum limit / target-side correspondence.}
  No PDE identification; no metric, spectral, Hecke, or
  cohomology correspondence theorem.
\item \textbf{Selection mechanism in the strong sense.}
  ``Cascade closure as selection mechanism'' is overstatement.
  This paper provides only a top-rung $(O1)$ witness for a
  chosen $\xi^{*}$ on the $H_4$ branch at finite $N$.
\end{itemize}

% =====================================================================
\section{Named open hypotheses}
\label{sec:open-hyp}
% =====================================================================

\begin{remark}[\textbf{(H-uniqueness)}: single-ray selection]
\label{rmk:H-uniqueness}
Selection of a single ray in $(F^0)^{W(H_4)} \setminus \{0\}$
requires structure beyond the cascade-closure flow. Candidates
for an additional selection principle include:
\begin{itemize}
\item closure-energy normalisation on the kernel of
  $L_N^{\mathrm{block}}$;
\item an additional symmetry beyond $W(H_4)$ (e.g.,
  time-reversal, or the $\sigma$ Galois action on rationality
  data of \cite[Definition~\texttt{def:sigma}]{RefinementSpaces});
\item a top-rung initial condition tied to a physical observable;
\item a thermodynamic / variational principle on the family of
  admissible $\xi^{*}$.
\end{itemize}
\textbf{(H-uniqueness)} is named here, not discharged. Per
\cite[\S 3.4 lines 637--639]{CascadeMechanism}, ``selection of
$\phi^{\star}_{N}$ itself is not provided by cascade closure and
is the successor paper's own task'' --- this paper supplies
\emph{a} $\phi^{\star}_{N}$ (parametrised by $\xi^{*}$), not
the unique one. Reducing the $28$-dimensional freedom to a single ray
lies \emph{outside} the cascade-closure consumer API.
\end{remark}

\begin{remark}[\textbf{(H-physical-interpretation)}]
\label{rmk:H-phys}
Whether $\iota_N\bigl((F^0)^{W(H_4)}\bigr) \subset X_N^{0, H_4}$
is the \emph{physically correct} selected subspace --- as opposed
to a mathematical artefact of restricting to the $H_4$ branch
with the constant-function embedding --- is a physical-modelling
question outside this paper's scope. The mathematical theorem
makes no claim of physical correctness or uniqueness of
interpretation.
\end{remark}

% =====================================================================
\section{Citation contract and downstream handoff}
\label{sec:handoff}
% =====================================================================

\subsection{What downstream papers may cite}

\begin{enumerate}[leftmargin=2em]
\item Definition~\ref{def:block-laplacian} ($L_N^{\mathrm{block}}$
  as an auxiliary operator on $X_N^{0, H_4}$, distinct from
  upstream operators by Remark~\ref{rmk:no-identification}).
\item Proposition~\ref{prop:psd} (PSD self-adjointness of
  $L_N^{\mathrm{block}}$ with explicit norm bound).
\item Theorem~\ref{thm:kernel} ($\ker L_N^{\mathrm{block}} \cong
  F^0$ via $\iota_N$).
\item Theorem~\ref{thm:convergence} (convergence of
  $e^{-t L_N^{\mathrm{block}}}$ with rate
  $\lambda_{+}(L_N^{\mathrm{block}})$).
\item Theorem~\ref{thm:fixed-dim} (explicit
  $\dim_{\R} (F^0)^{W(H_4)} = 28$).
\item Theorem~\ref{thm:main} (witnessed $(O1)$ discharge for
  any chosen $\xi^{*} \in (F^0)^{W(H_4)} \setminus \{0\}$).
\item Worked numerics of \S\ref{sec:numerics} ($N = 0$ explicit
  values).
\end{enumerate}

\subsection{What downstream papers may not cite}

\begin{enumerate}[leftmargin=2em]
\item Any theorem about $L_N^{\mathrm{block}}$ as a theorem about
  P4's vertex / top-left block
  $A_n^{\mathrm{vertex}} = d_n^{*} d_n$ or compat's
  $\widetilde A_n$: the operators differ
  (Remark~\ref{rmk:no-identification}).
\item Single-state selection of a unique
  $\phi^{\star}_{N}$: Theorem~\ref{thm:main} parametrises
  $\phi^{\star}_{N}$ by $\xi^{*} \in (F^0)^{W(H_4)} \setminus
  \{0\}$ (a $28$-dimensional cone); single-ray selection is
  \textbf{(H-uniqueness)} (\S\ref{sec:open-hyp}).
\item Inverse-limit $(O1)$ as $N \to \infty$: open under
  $(L1), (L3)$
  (\cite[\S\texttt{sec:scope-and-gap}]{RefinementCompat}).
\item $(O2), (O3)$, the full cascade-closure $(O0)$--$(O3)$ API:
  not addressed (see Remark~\ref{rmk:non-claims}).
\item $D_4$-branch witness: requires a separate paper.
\end{enumerate}

\subsection{Recommended citation phrasing}

``By cascade-selection-h4 (Theorem~\ref{thm:main}), the
cascade-mechanism consumer-API obligation $(O1)$ on the $H_4$
branch is discharged at every finite top-rung level $N \geq 0$ by
the explicit five-tuple
$(X_N^{0, H_4}, \tfrac12 \langle \cdot, L_N^{\mathrm{block}}
\cdot \rangle, e^{-t L_N^{\mathrm{block}}}, \iota_N(\xi^{*}),
\iota_N(\xi^{*}))$ for any chosen $\xi^{*} \in (F^0)^{W(H_4)}
\setminus \{0\}$, where $L_N^{\mathrm{block}}$ is the auxiliary
block graph Laplacian and $(F^0)^{W(H_4)}$ has explicit real
dimension $28$. The construction leaves a $28$-dimensional free
choice of $\xi^{*}$; single-ray selection is the named open
hypothesis \textbf{(H-uniqueness)}, outside the cascade-closure
consumer API.''

% =====================================================================
\begin{thebibliography}{99}

\bibitem{RefinementSpaces}
\emph{Cascade--Classical Correspondence III: Refinement Spaces and
Symmetry Actions}, paper P3.
\texttt{papers/cascade-refinement-spaces/cascade-refinement-spaces.tex}.
Pinned results used here, cited by label:
\texttt{def:Xn-0}, \texttt{def:refinement}, \texttt{def:base-graphs},
\texttt{prop:Xn-Hilbert}, \texttt{not:fibres},
\texttt{def:coxeter}, \texttt{lem:coxeter-unitary},
\texttt{def:p0}, \texttt{thm:bonding}, \texttt{def:sigma}.

\bibitem{ClosureDynamics}
\emph{Cascade--Classical Correspondence IV: Discrete Closure
Dynamics}, paper P4.
\texttt{papers/cascade-closure-dynamics/cascade-closure-dynamics.tex}.
Pinned results used here, cited by label:
\texttt{def:coboundary} (the $F^0 \to F^1 = \mathrm{Im}(\OO)$
projection), \texttt{def:An} (the vertex / top-left block
$A_n^{\mathrm{vertex}} = d_n^{*} d_n$),
\texttt{def:Fn} (P4's closure functional), and
\texttt{thm:flow-exists} (style-only: same finite-dim ODE /
matrix-exponential argument used in
Proposition~\ref{prop:o0}(O0.3) for $L_N^{\mathrm{block}}$).

\bibitem{RefinementCompat}
\emph{The Cascade Curvature Operator under Refinement:
Schur-Complement Identity and Energy Intertwining}.
\texttt{papers/cascade-refinement-compat/cascade-refinement-compat.tex}.
Pinned results used here, cited by label:
\texttt{def:abstract-model} (the abstract scalar pure-midpoint
refinement model), \texttt{def:Atilde} (the scaled curvature
operator $\widetilde A_n = 2^n A_n^{\mathrm{vertex}}$ in that
abstract model), \texttt{prop:strict-fails} ($(O2)$ as stated
fails for nontrivial refinements),
\texttt{thm:O3-discharge} ($(O3)$ closed in the abstract scalar
model), \texttt{sec:scope-and-gap} (the open lift hypotheses
$(L1)$ and $(L3)$).

\bibitem{CascadeMechanism}
\emph{Cascade Closure: A Conditional Compatibility Template for
Geometric State Selection}, paper P-mech.
\texttt{papers/cascade-mechanism/cascade-mechanism.tex}.
Pinned results used here: the consumer API $(O0)$--$(O3)$
(\S 3.4 lines 594--627), with $(O1)$ specifically at lines
609--617, and the explicit ``selection is the successor's task''
framing at lines 637--639.

\bibitem{AlgebraicSubstrate}
\emph{Cascade--Classical Correspondence II: The Cascade Algebraic
Substrate}, paper P2.
\texttt{papers/cascade-algebraic-substrate/cascade-algebraic-substrate.tex}.
Pinned results used here, cited by label:
\texttt{def:V600} ($V_{600}$ vertex set);
\texttt{def:icosian-ring} (the icosian ring $I \subset \HH$);
\texttt{prop:24-in-600} ($V_{24} \subset V_{600}$);
\texttt{tab:P2-spectrum} (the $V_{600}$ scalar Laplacian
spectrum, used only in the $N = 0$ worked numerics of
\S\ref{sec:numerics}). \textbf{No P2 character data is
load-bearing for the main theorem}; irreducibility of the
$H_4$-reflection representation is taken from \cite{Humphreys}.

\bibitem{Humphreys}
J.~E.~Humphreys, \emph{Reflection Groups and Coxeter Groups},
Cambridge Studies in Advanced Mathematics \textbf{29}, Cambridge
University Press, 1990.
$H_4$ as a finite irreducible Coxeter group: Ch.~2;
irreducibility of the standard reflection representation: Ch.~5.

\bibitem{Chung}
F.~R.~K.~Chung, \emph{Spectral Graph Theory}, CBMS Regional
Conference Series in Mathematics 92, AMS, 1997.
Combinatorial graph Laplacian, kernel structure, and
operator-norm bound on (multi)graphs: Ch.~1.

\bibitem{ReedSimonI}
M.~Reed and B.~Simon, \emph{Methods of Modern Mathematical
Physics, Vol.~I: Functional Analysis}, rev.\ ed., Academic
Press, 1980.
Spectral theorem and functional calculus for self-adjoint
finite-dim operators: Ch.~VIII.5; bounded operators on Hilbert
spaces: Ch.~II.

\bibitem{HornJohnson}
R.~A.~Horn and C.~R.~Johnson, \emph{Matrix Analysis}, 2nd ed.,
Cambridge University Press, 2013.
Matrix exponentials and linear ODEs: Ch.~6.2; spectrum and
operator norm of block-diagonal matrices: Ch.~6.4.

\bibitem{CoxeterRegularPolytopes}
H.~S.~M.~Coxeter, \emph{Regular Polytopes}, 3rd ed., Dover, 1973.
$600$-cell edge graph: \S 22.7;
$24$-cell: \S 22.2;
Schl\"{a}fli refinement: \S 8.

\end{thebibliography}

\end{document}
